% =============================================================================
% Copyright (c) 2026 Mert Torun, M.Sc. (mtorun0x7cd) <info@mtorun0x7cd.com>
% 
% Permission is hereby granted, free of charge, to any person obtaining a copy
% of this software and associated documentation files (the "Software"), to deal
% in the Software without restriction, including without limitation the rights
% to use, copy, modify, merge, publish, distribute, sublicense, and/or sell
% copies of the Software, and to permit persons to whom the Software is
% furnished to do so, subject to the following conditions:
% 
% The above copyright notice and this permission notice shall be included in
% all copies or substantial portions of the Software.
% 
% THE SOFTWARE IS PROVIDED "AS IS", WITHOUT WARRANTY OF ANY KIND, EXPRESS OR
% IMPLIED, INCLUDING BUT NOT LIMITED TO THE WARRANTIES OF MERCHANTABILITY,
% FITNESS FOR A PARTICULAR PURPOSE AND NONINFRINGEMENT. IN NO EVENT SHALL THE
% AUTHORS OR COPYRIGHT HOLDERS BE LIABLE FOR ANY CLAIM, DAMAGES OR OTHER
% LIABILITY, WHETHER IN AN ACTION OF CONTRACT, TORT OR OTHERWISE, ARISING FROM,
% OUT OF OR IN CONNECTION WITH THE SOFTWARE OR THE USE OR OTHER DEALINGS IN THE
% SOFTWARE.
% =============================================================================

\chapter{Cross-Platform Environment Manual}
\label{app:manual}

This appendix provides a step-by-step setup and configuration guide for running the \textit{OneWare Container Extension} across different platforms and container runtimes.

\section{Installation}
To install the Container Extension within OneWare Studio, execute the following steps:
\begin{enumerate}
    \item Launch OneWare Studio.
    \item Navigate to the \textbf{Extension Manager} panel.
    \item Search for the ``Container Extension'' in the marketplace and click \textbf{Install} to download and install it directly.
    \item Restart the IDE to load the background daemon services and registration hooks.
\end{enumerate}

\section{Global Configuration Settings}
The Container Extension allows users to customize container parameters and security bounds. Table~\ref{tab:config_keys} outlines all the settings registered in the OneWare configuration database by \texttt{ContainerExtensionModule.cs}.

\begin{table}[htbp]
    \caption{Configuration settings keys defined in \texttt{ContainerExtensionModule.cs}}
    \label{tab:config_keys}
    \centering
    \scriptsize
    \begin{tabularx}{\textwidth}{@{} p{4.5cm} p{2.5cm} X @{}}
        \toprule
        \textbf{Setting Key} & \textbf{Default} & \textbf{Description} \\ \midrule
        \path{ContainerExtension_DockerRuntimePath} & \texttt{""} & Absolute path to the runtime executable (defaults to host \texttt{PATH}). \\
        \path{ContainerExtension_DefaultImage} & \path{"fentwums/oss-cad-suite:latest"} & The default full-flow oss-cad-suite image for all tools; build-only (not on Docker Hub), produced via Build Local Image rather than pulled. \\
        \path{ContainerExtension_MemoryLimit} & \texttt{0.0} & RAM limit in MB (0 = unlimited). \\
        \path{ContainerExtension_CpuLimit} & \texttt{0.0} & CPU cores limit (0 = unlimited). \\
        \path{ContainerExtension_Timeout} & \texttt{0.0} & Maximum container execution runtime in minutes (0 = no limit). \\
        \path{ContainerExtension_NetworkMode} & \texttt{"bridge"} & Network driver mode (\texttt{bridge}, \texttt{host}, or \texttt{none}). \\
        \path{ContainerExtension_PullPolicy} & \texttt{"if-\allowbreak not-\allowbreak present"} & Image pull policy (\texttt{always}, \texttt{if-not-present}, or \texttt{never}). \\
        \path{ContainerExtension_AutoRemove} & \texttt{true} & Automatically remove container after execution (\texttt{--rm}). \\
        \path{ContainerExtension_AllowPrivileged} & \texttt{false} & Allow privileged container execution. \\
        \path{ContainerExtension_LogLevel} & \texttt{"Errors Only"} & Telemetry verbosity (\texttt{Off}, \texttt{Errors Only}, \texttt{Info}, \texttt{Verbose}). \\
        \path{ContainerExtension_ShowTimestamps} & \texttt{true} & Prepend timestamp to stdout/stderr stream. \\
        \path{ContainerExtension_ContainerNamePrefix} & \texttt{"container\-extension-"} & Prefix for spawned container names. \\
        \path{ContainerExtension_ExtraFlags} & \texttt{""} & Custom runtime flags or labels to append. \\
        \path{ContainerExtension_DashboardRefresh} & \texttt{"Manual"} & Dashboard statistics polling interval. \\
        \path{ContainerExtension_TelemetryRetention} & \texttt{"25"} & Number of recent executions retained in telemetry cache. \\
        \path{ContainerExtension_DaemonSocket} & \texttt{""} & Custom socket endpoint override (e.g., \texttt{unix://...} or \texttt{npipe://...}). \\
        \path{ContainerExtension_BypassNamedPipeCheck} & \texttt{true} & Bypass the Windows named-pipe server-trust validation. Default-on: the Windows-only check yields false positives on WSL2 relays and rootless/remote engines; uncheck it on a hardened Windows host to re-enable the guard. \\
        \path{ContainerExtension_AllowNativeFallback} & \texttt{false} & Fall back to native host execution if container daemon is offline. \\
        \path{ContainerExtension_Platform} & \texttt{"auto"} & Overrides container architecture target (e.g., \texttt{linux/amd64}). \\
        \bottomrule
    \end{tabularx}
\end{table}

\newpage
\section{Runtime-Specific Configuration Manuals}

\subsection{OrbStack (macOS)}
OrbStack is a highly optimized, lightweight Docker Desktop alternative for macOS. It integrates natively with the macOS Hypervisor framework, offering superior I/O performance via VirtioFS.

\begin{enumerate}
    \item \textbf{Socket Mapping:} OrbStack establishes a standard Unix domain socket at \path{/var/run/docker.sock}, which is automatically symlinked. If OneWare cannot locate the socket, set the configuration key:
    \begin{lstlisting}[language=bash]
ContainerExtension_DaemonSocket = unix:///var/run/docker.sock
\end{lstlisting}
    \item \textbf{Architecture Emulation (Apple Silicon):}
    When running on Apple Silicon (M-Series chip), running AMD64 images incurs a translation tax via Rosetta 2. For maximum performance, configure the platform override to native ARM64:
    \begin{lstlisting}[language=bash]
ContainerExtension_Platform = linux/arm64
\end{lstlisting}
    Alternatively, keep it on \texttt{auto} to allow the extension to query the daemon and determine availability.
    \item \textbf{Resource Optimization:} Ensure OrbStack's shared memory limit matches or exceeds the workspace requirements. If a synthesis run fails due to insufficient memory, set:
    \begin{lstlisting}[language=bash]
ContainerExtension_MemoryLimit = 8192 # Limits container to 8 GB
\end{lstlisting}
\end{enumerate}

\subsection{Docker Desktop (Windows and macOS)}
Docker Desktop is the official cross-platform container management system.

\begin{enumerate}
    \item \textbf{Named Pipe Impersonation (Windows):}
    On Windows hosts, Docker Desktop communicates via named pipes, typically at \texttt{npipe://./pipe/docker\_engine}. The client always opens the pipe at \path{TokenImpersonationLevel.Identification}, so the daemon endpoint can authenticate the caller but cannot impersonate its token. On top of this, an optional server-trust check verifies the pipe server process before the connection is used.
    \item \textbf{Bypassing the Server-Trust Check:}
    That server-trust check is Windows-only and produces false positives on common configurations (WSL2 relays, rootless or remote engines), so it is \emph{bypassed by default} (\path{ContainerExtension_BypassNamedPipeCheck = true}). On a hardened Windows host where the daemon is local and trusted, uncheck ``Bypass Named Pipe Security Check'' to re-enable it:
    \begin{lstlisting}[language=bash]
ContainerExtension_BypassNamedPipeCheck = false
\end{lstlisting}
    \item \textbf{WSL 2 Integration:}
    For Windows hosts, ensure WSL 2 is enabled in the Docker Desktop settings and that the target Linux distribution is ticked under \textbf{Resources -> WSL Integration} to ensure correct absolute directory path mappings.
\end{enumerate}

\subsection{Podman (Rootless Linux Setup)}
Podman is a daemonless, rootless container engine that operates natively on Linux.

\begin{enumerate}
    \item \textbf{Rootless Socket Setup:}
    To run Podman without root privileges, enable the user-level systemd service:
    \begin{lstlisting}[language=bash]
systemctl --user enable --now podman.socket
\end{lstlisting}
    This exposes a local user socket at \texttt{/run/user/\$UID/podman/podman.sock}.
    \item \textbf{Daemon Connection Config:}
    Configure the OneWare settings to point to the user-space socket:
    \begin{lstlisting}[language=bash]
ContainerExtension_DaemonSocket = unix:///run/user/1000/podman/podman.sock
ContainerExtension_DockerRuntimePath = /usr/bin/podman
\end{lstlisting}
    Replace \texttt{1000} with the target host user ID.
    \item \textbf{User Namespace Mapping:}
    Rootless execution requires correct UID mapping. Verify that \texttt{/etc/subuid} and \texttt{/etc/subgid} have ranges defined for the execution user:
    \begin{lstlisting}[language=bash]
user:100000:65536
\end{lstlisting}
    This mapping translates host user ID \texttt{1000} inside the container boundary seamlessly, resolving host-to-container write permission issues on bound output directories.
\end{enumerate}

\section{Troubleshooting Guide}

\subsection{Socket Permissions (Linux)}
If compilation fails with socket permission errors, ensure the current user belongs to the \texttt{docker} group:
\begin{lstlisting}[language=bash]
sudo usermod -aG docker $USER
newgrp docker
\end{lstlisting}
If the group configuration is restricted by system administrators, apply temporary read-write permissions to the socket file:
\begin{lstlisting}[language=bash]
sudo chmod 666 /var/run/docker.sock
\end{lstlisting}

\subsection{Dangling Container Reap}
On start, the Container Extension scans for and reaps any orphaned containers matching the user-configured name prefix. If dangling containers are leaking memory or host resources, verify the prefix:
\begin{lstlisting}[language=bash]
ContainerExtension_ContainerNamePrefix = containerextension-
\end{lstlisting}
Spawned containers are automatically reaped after run completion unless the configuration key \path{ContainerExtension_AutoRemove} is set to \texttt{false}.
